% ===========================================================================
%  CH10/main.tex --- the conclusion chapter, written in full in this file.
%
%  Everything between the CHAPTER-BODY markers below is the thesis-ready
%  fragment: it opens with \chapter and uses house notation only.  The copy
%  under Full_stitch_1/chapters/ch10/chapter.tex is generated from those
%  markers and should not be edited by hand.
%
%  Cross-chapter references are written with \mextref{label}{fallback}, so
%  that this standalone compile prints the fallback text and the stitched
%  build resolves the reference.
% ===========================================================================

\documentclass[11pt,a4paper]{report}

\usepackage[a4paper,margin=1in]{geometry}
\usepackage[T1]{fontenc}
\usepackage[utf8]{inputenc}
\usepackage{lmodern}
\usepackage{microtype}
\usepackage{amsmath,amssymb,amsthm,mathtools}
\usepackage{graphicx}
\usepackage{float}
\usepackage{placeins}
\usepackage{booktabs}
\usepackage{array}
\usepackage[hypcap=false]{caption}
\usepackage{tikz}
\usetikzlibrary{calc,arrows.meta,positioning}
\usepackage{enumitem}
\usepackage{xcolor}
\usepackage{csquotes}
\usepackage{hyperref}
\usepackage[nameinlink,noabbrev]{cleveref}

\setcounter{secnumdepth}{3}
\setcounter{tocdepth}{2}
\allowdisplaybreaks
\setlength{\parindent}{1.2em}
\setlength{\parskip}{0.15em}

\crefformat{equation}{(#2#1#3)}
\Crefformat{equation}{(#2#1#3)}

\definecolor{ink}{HTML}{1A1C1F}
\definecolor{inksoft}{HTML}{565B62}
\definecolor{accent}{HTML}{0072B2}
\definecolor{type2}{HTML}{D55E00}

% --- House notation (matches the stitched preamble) ------------------------
\newcommand{\pr}[1]{\mathbb{P}\!\left(#1\right)}
\newcommand{\ex}[1]{\mathbb{E}\!\left(#1\right)}
\newcommand{\mean}[1]{\mathbb{E}\!\left[#1\right]}
\newcommand{\var}[1]{\operatorname{Var}\!\left(#1\right)}
\newcommand{\dd}{\mathrm{d}}
\newcommand{\ee}{\mathrm{e}}
\newcommand{\lrb}[1]{\left(#1\right)}
\newcommand{\lrbs}[1]{\left[#1\right]}
\newcommand{\Kb}{\mathcal{K}}
\newcommand{\Ifix}[1][]{I_{\mathrm{fix}#1}}
\newcommand{\ft}{\textit{F.~tularensis}}
\newcommand{\yp}{\textit{Y.~pestis}}
\providecommand{\mextref}[2]{%
  \ifcsname r@#1\endcsname\cref{#1}\else#2\fi}
\providecommand{\Mextref}[2]{%
  \ifcsname r@#1\endcsname\Cref{#1}\else#2\fi}

\theoremstyle{plain}
\newtheorem{theorem}{Theorem}[chapter]
\newtheorem{proposition}[theorem]{Proposition}
\theoremstyle{remark}
\newtheorem{remark}[theorem]{Remark}

\hypersetup{
  colorlinks=true,
  linkcolor=blue!55!black,
  citecolor=red!55!black,
  urlcolor=blue!60!black,
  pdftitle={Conclusion -- the evolved strategy of Yersinia pestis},
  pdfauthor={Adam Aldridge}
}

\begin{document}

\hypersetup{pageanchor=false}
\begin{titlepage}
\centering
\vspace*{4cm}
{\LARGE\bfseries Conclusion\par}
\vspace{6pt}
{\large The evolved strategy of \textit{Yersinia pestis}\par}
\vspace{2.5cm}
{\large Adam Aldridge\par}
\vspace{4pt}
{University of Leeds\par}
\end{titlepage}
\hypersetup{pageanchor=true}

\tableofcontents
\clearpage

\setcounter{chapter}{9}   % standalone prints as Chapter 10

% ===========================================================================
%<<<CHAPTER BODY START>>>
\chapter[Conclusion]{Conclusion: the evolved strategy of
\texorpdfstring{\textit{Yersinia pestis}}{Yersinia pestis}}
\label{ch:conclusion}

This thesis began with one organism and one peculiarity. \yp{} enters a
mammalian host from a flea bite, permits itself to be taken up by
macrophages, replicates inside them for a while, and only then reverses
course and resists phagocytosis altogether. The reversal is not incidental
book-keeping. It is among the sharpest differences between \yp{} and the
comparatively mild ancestor it descends from, and it is the sort of thing
that invites a mathematical answer, because it is a choice between two
strategies whose returns can in principle be written down and compared. I
wanted to know whether the switch could be shown to pay, and to say in what
currency it pays.

Most of the thesis is the machinery that question turned out to need. The
choice is made by a small number of bacteria inside a single cell, where
whether the lineage survives at all is settled by chance rather than by a
mean, so the object required is a distribution and not a trajectory. Getting
that distribution for a cell that grows its occupants and then ruptures took
\mextref{bdc:ch:core}{Chapter~4} and then
\mextref{dist:ch:distribution-theory}{Chapter~5}; carrying it up to a
population took \mextref{p:ch:one-cell-to-population}{Chapter~6}; and asking
what the release schedule is worth against defences that can be exhausted
took \mextref{path:ch:pathogenesis}{Chapter~8}. Along the way two questions
detained the work longer than they strictly deserved, and both became
chapters of their own: a constant governing subcritical branching
conditioned on survival (\mextref{gw:ch:galton-watson}{Chapter~3}), and the problem of writing down
Markov chains whose rates are allowed to move
(\mextref{var:ch:rates}{Chapter~9}).

What follows is a reckoning rather than a summary. \Cref{con:sec:contributions}
sets out what each chapter established and what it left standing.
\Cref{con:sec:questions} returns to the four questions posed in
\mextref{ch:introduction}{Chapter~1} and answers them as far as the work
allows, which in two cases is not all the way.
\Cref{con:sec:quadrant} gives the one argument that never found a home in a
chapter of its own: a small piece of reasoning about why an organism might
want to be eaten early and not late. \Cref{con:sec:limits} says what the
thesis does not establish, and \cref{con:sec:closing} closes.

\section{Contributions, chapter by chapter}
\label{con:sec:contributions}

\Mextref{m:ch:math-intro}{Chapter~2} assembles the objects and the methods:
Markov chains in discrete and continuous time, the branching processes that
matter here, and the generating-function machinery --- forward and backward
equations, extinction probabilities, limiting conditional means, and the
method of characteristics --- each developed once, on whichever process shows
it most plainly. Two of its appendices are more than review. The first
records how heavily survival near criticality depends on the size of the
founding cohort (\mextref{m:app:smallpops}{Section~2.A}), which is the fact
that almost everything downstream leans on. The second solves a family of
absorption models in which particles are drawn into a compartment and may
then replicate inside it; the hardest member of that family is solved here
for the first time (\mextref{m:app:absorption}{Section~2.B}).

\Mextref{gw:ch:galton-watson}{Chapter~3} is the chapter furthest from plague,
and it is kept because the question it answers is a real one and the answer
is complete. For the subcritical binary Galton--Watson process the
survival probability decays as $S_n\sim A(p)(2p)^n$, and that single constant
also fixes the population conditioned on survival, whose mean converges to
$1/A(p)$. The constant exists because the limit can be rewritten as a
convergent infinite product. Rescaling the survival recursion to the logistic
map identifies it with a value of the associated Koenigs linearising
function, and a theorem of Becker and Bergweiler shows that function to be
hypertranscendental (\mextref{gw:thm:main}{Theorem~3.5}), which does not
settle the closed-form question for $p\mapsto A(p)$ but does explain why a
long search for a formula came to nothing. What remains is practical: the
product, together with the near-critical asymptotic $A(p)\sim2(1-2p)$,
describes the constant across the whole subcritical range.

\Mextref{bdc:ch:core}{Chapter~4} introduces the process the rest of the
thesis is built on. Inside one host cell each pathogen particle divides at
rate $\lambda$, dies at rate $\mu$, and triggers rupture of the whole cell at
rate $\delta$. The calculation closes because the probability of no rupture
yet, $I(t)$, obeys an autonomous Riccati equation; the two roots $a$ and $b$
of its characteristic quadratic then carry every closed form in the thesis.
Cell-fate probabilities, the mean and second moment of the intracellular
load, the mean release and the first two conditional moments of the burst
size all follow in that one coordinate, and several of them --- the second
moment expressed through $I$, the variances, the joint load--release moments
--- go beyond what the birth--death--catastrophe literature had recorded.

\Mextref{dist:ch:distribution-theory}{Chapter~5} completes the single cell in
distribution. The load is geometric at every time, with a ratio that slides
towards $1/a$; the quasi-stationary law is that endpoint
(\mextref{dist:thm:qsd}{Theorem~5.6}); and the burst size, averaged over the
instant of rupture, is the same law
(\mextref{dist:thm:identity}{Theorem~5.10}). That the two coincide is the most
satisfying result in the thesis, because the two constructions share no step:
one conditions a surviving trajectory on never being absorbed, the other
averages over rupture times, and the identity turns out to hold because
size-biasing telescopes the occupation integral onto the endpoint of the
process's own geometric ratio. Two probes then mark the boundary of the
structure. Several founders sharing a cell leave it intact; chaining one
cell's burst into the next destroys it as soon as internal extinction is
possible, and that case remains open.

\Mextref{p:ch:one-cell-to-population}{Chapter~6} embeds the cell in a
population. The survival and release statistics become age-structured
kernels, and the resulting renewal system is exact in the mean. Three
consequences follow. The standard model of within-host virus dynamics is a
projection of that system, valid one exponential phase at a time
(\mextref{p:thm:projection}{Theorem~6.3}), so a fitted release rate is a
property of the phase and not of the organism, and a model that already
carries a hand-fitted eclipse alongside a bursting pathogen is fitting the
same delay twice. No fit of the two population constants recovers the three
intracellular rates (\mextref{p:sec:identifiability}{\S6.5.6}), so a
population fit is worth doing only alongside a single-cell measurement.
And at fixed total output under linear clearance the release schedule is
invisible to $R_0$ (\mextref{p:thm:R0inv}{Theorem~6.6}) while remaining
visible to establishment: bursting strictly lowers the extinction
probability of a small inoculum if and only if the flooding parameter
$L=a(1-b)$ exceeds one (\mextref{p:thm:flood}{Theorem~6.8}).

\Mextref{tt:ch:two-type}{Chapter~7} solves a two-type version exactly. Type~1
converts irreversibly into type~2, either type may trigger a single global
catastrophe, and because the catastrophe rate depends on the whole path the
natural object is an occupation-time transform rather than a terminal-count
generating function. The backward system is triangular, and after
linearisation the surviving Riccati equation becomes the Gauss hypergeometric
equation, giving a closed formula for the finite-time probability of avoiding
catastrophe (\mextref{tt:thm:main}{Theorem~7.4}) which holds for any ordering
of the type-specific rates and reduces to the two-type branching equations of
Antal and Krapivsky \cite{antal2011exact} when the catastrophe rates vanish.
Read biologically, the two types are early and adapted intracellular
phenotypes and the catastrophe is irreversible loss of containment.

\Mextref{path:ch:pathogenesis}{Chapter~8} asks what a burst is worth when the
defence can be used up. Inside the cell, a replicator--suppressor process
with a granule store that is never replenished makes a larger founding cohort
disproportionately safer, because the cohort itself spends the store. Outside
it, a fixed removal budget $\vartheta$ per delivery gives the surviving
release in closed form, $\mathcal{R}(\vartheta)=V_\infty a^{-\vartheta}$, and
with it a reversal: against a threshold rather than a
rate, the flooding criterion flips, and bursting beats matched-mean budding
precisely when $L<1$ (\mextref{path:prop:reversal}{Proposition~8.12}). The
reason is a single fact about variance. Extinction probability rewards low
offspring variance whereas the surplus above a threshold is convex and
rewards high variance, so the sign of $L-1$ drives both criteria and drives
them opposite ways. The chapter also yields the number that reappears below:
a removal budget $\vartheta^{\ast}=\log V_\infty/\log a$ beyond which no
intracellular strategy stays supercritical
(\mextref{path:cor:threshold}{Corollary~8.8}).

\Mextref{var:ch:rates}{Chapter~9} lets the rates move. Under a
time-varying per-capita birth rate the species count stays geometric and
depends on the rate law only through its integral, which under logistic
growth converges and leaves a proper limit law rather than merely a finite
late-time mean (\mextref{var:cor:logistic-limit}{Corollary~9.3}). In a model
of Pimentel's competition experiment the state at which the chain is equally
likely to step up as down is available in closed form and is unstable
(\mextref{var:eq:icrit}{(9.21)}), which is a different thing from a process
approaching an equilibrium. Two further models are carried only as far as
simulation. The chapter's general observation is worth repeating, because it
is not the obvious one: what costs the generating function is not feedback as
such but \emph{continuous bending} of a rate. A process may act on its own
circumstances as much as it likes and stay solvable, provided what survives
the interaction is a random time or a random count rather than a moving rate.

\begin{table}[!htbp]
\centering
\footnotesize
\renewcommand{\arraystretch}{1.15}
\begin{tabular}{@{}l>{\raggedright\arraybackslash}p{0.40\textwidth}>{\raggedright\arraybackslash}p{0.36\textwidth}@{}}
\toprule
Ch. & What it establishes & What it leaves \\
\midrule
2 & Objects and methods; survival of founding cohorts near criticality;
    a family of absorption models solved
  & The absorption models are not used by the modelling chapters \\
3 & $A(p)$ exists, equals a Koenigs value, has no elementary formula in $z$,
    and is computable across $(0,\tfrac12)$
  & Whether $p\mapsto A(p)$ has a closed form \\
4 & The birth--death--catastrophe cell: fates, moments, mean release, burst
    moments, all through $I(t)$
  & Rates unmeasured for \yp{}; no data confrontation \\
5 & Geometric load at every time; quasi-stationary law; burst-size law
    identical to it, with a structural proof
  & The chained process for $\mu>0$; whether the telescoping criterion is
    necessary \\
6 & An exact burst-aware renewal model; the standard model as its
    projection; non-identifiability; the flooding criterion $L>1$
  & Where \yp{} sits relative to $L=1$; multiplicity through the population
    layer \\
7 & Exact hypergeometric no-catastrophe probability for two types with one
    global catastrophe
  & Killed second moments, which an eclipse-aware renewal model needs \\
8 & Depletable defences: certainty threshold, $\mathcal{R}(\vartheta)$, the
    reversal at $L=1$, and the budget $\vartheta^{\ast}$
  & Population-level consequences of a threshold clearance \\
9 & Distribution under a moving birth rate; logistic limit law; an unstable
    moving critical point
  & Four of the models are specifications or simulations only \\
\bottomrule
\end{tabular}
\caption{What each chapter settles, and what it hands on. The right-hand
column is drawn from the chapters' own statements of scope, not added here.}
\label{con:tab:contributions}
\end{table}

\FloatBarrier

\section{The four questions, answered}
\label{con:sec:questions}

\subsection{Why \texorpdfstring{\yp{}}{Y. pestis} switches phase}
\label{con:sec:q1}

Two arguments in this thesis reach the switch, by different routes, and they
agree on its shape if not yet on its position. The first is the quadrant
argument of \cref{con:sec:quadrant}: with three environments ordered by
hostility, the expected return on accepting phagocytosis falls linearly as
the neutrophil fraction rises, crosses the constant return on resisting it,
and the crossing sits strictly inside $(0,1)$. The second is the removal
budget of \mextref{path:sec:ypestis}{\S8.5.4}: while the killer-cell density
is roughly constant a cell releasing its whole load at once pays one toll
where a budding cell would pay many, and once a rise in that density pushes
the toll past $\vartheta^{\ast}=\log V_\infty/\log a$ no intracellular
strategy remains supercritical. A third strand supports both without
deciding either. Just above criticality the size of a founding cohort weighs
very heavily on whether a lineage survives at all
(\mextref{m:app:smallpops}{Section~2.A}), and the intracellular phase is
exactly a device for arriving in numbers.

The common prediction is worth isolating, because it is the part that could
be wrong in an interesting way. Both arguments key the switch to \emph{how
much of the killer field is present}, not to a clock: to the fraction of
phagocytes that are neutrophils in one case, to the number of agents removed
per delivery in the other. Neither predicts a fixed interval after inoculation. The histology
does place the switch alongside the neutrophil response
\cite{BosioEtAl2005,PetersEtAl2013}, which is consistent with both accounts
and discriminates between neither, and whether the neutrophils drive the
switch or merely arrive with it is not settled by those studies.

Three things cannot be claimed. The two thresholds have not been shown to be
consistent --- whether the neutrophil fraction that crosses $p_N^{\ast}$ is
the one that carries the toll past $\vartheta^{\ast}$ is an open problem
recorded in \mextref{path:sec:open}{\S8.6.2} and repeated in
\cref{con:sec:limits}. Neither threshold has been put against data. And
neither argument excludes the plainer explanations: that macrophages are
being used as transport to the draining lymph nodes, or that the early
pro-phagocytic phase is a hangover from life inside the flea. What the
mathematics supplies is not a proof that the strategy is optimal but a
statement of what would have to be measured for the question to be settled.

\subsection{Whether bursting pays}
\label{con:sec:q2}

It does, but ``bursting pays'' turns out not to be one claim. Four
comparisons run through \mextref{p:ch:one-cell-to-population}{Chapter~6} and
\mextref{path:ch:pathogenesis}{Chapter~8}, and they do not agree with one
another, which is the finding rather than a defect.

At fixed total output per cell and linear clearance, replacing a constant
drip by a single burst leaves $R_0$ exactly unchanged: the schedule is
invisible to the reproduction number. At the same fixed mean, the schedule is
nevertheless not invisible to establishment, and bursting strictly lowers the
extinction probability of a small inoculum precisely when $L=a(1-b)$ exceeds
one. Against a defence with a fixed removal budget rather than a clearance
rate, that criterion reverses exactly, and bursting wins when $L<1$. A fourth
comparison, holding the release law fixed and asking how the same total fares
when split across more deliveries, favours concentration unconditionally,
because a toll paid per delivery is paid once by a bursting cell and $n$
times by a budding one.

So the honest answer is that a depletable defence supplies a genuine
selection gradient towards fewer and larger deliveries, and that whether one
particular bursting law beats one particular budding law at matched mean
depends on whether the defence is a rate or a budget, with $L$ the quantity
that decides. What remains out of reach is placing \yp{} on that scale. Doing
so needs $\lambda$, $\mu$ and $\delta$ for the macrophage phase, and $\delta$
is fixed
by the exit mechanism, which is exactly the thing the primary literature
still treats as unresolved --- apoptosis, pyroptosis and necroptosis are all
reported \cite{MonackEtAl1997,ZaubermanEtAl2006,PetersEtAl2013}. The two
parameter sets available in the right form belong to \ft{}
\cite{carruthers2020stochastic} and to \textit{B.~anthracis}
\cite{williams2021anthrax}, and they describe cells at opposite extremes of
what the model allows --- near-certain rupture releasing hundreds in the one
case, near-certain clearance with minimal release in the other. Neither is
\yp{}.

\subsection{An intracellular replacement for the standard model}
\label{con:sec:q3}

Yes, and it is more than a substitution of one rate for another. The standard
model gives an infected cell two constants, a release rate and a removal
rate, and for a budding virus that is an exact description rather than an
approximation. For a lytic pathogen it is wrong in mechanism, because release
and death are not independent events but the same event. The renewal system
of \mextref{p:ch:one-cell-to-population}{Chapter~6} replaces the two
constants with two kernels derived from the intracellular process, is exact
in the mean, and returns the standard model as a projection of itself, valid
one exponential phase at a time. Three consequences bear directly on fitting
practice: acute-phase and set-point estimates of the release rate should
differ systematically and by a bounded amount; the eclipse emerges from the
intracellular dynamics with no extra parameter, so a model carrying a fitted
one may be counting the same delay twice; and because the population fit is
not identifiable for the three intracellular rates, it is worth doing only
alongside a single-cell companion measurement. The single-cell predictions
that would supply that companion are listed in
\mextref{dist:sec:discussion-assay}{\S5.12.1}, and are sharp enough to be
wrong.

\subsection{Rates that move}
\label{con:sec:q4}

Partially, and the partial answer has a clean edge to it. Three results in
\mextref{var:ch:rates}{Chapter~9} carry closed forms --- the geometric law
under a moving birth rate, the logistic-speciation limit law, and the
unstable moving critical point of the Pimentel chain --- and the remaining
models are specifications or simulations that outrun the machinery of the
earlier chapters. What that collection shows is where the machinery gives
out, and the boundary is not the expected one. It is not feedback that costs
the generating function. A process may consume a
capacity, harden a competitor or infect its neighbours and still be solvable,
provided the interaction leaves behind a random time or a random count rather
than a rate that goes on bending. That is a modest observation, but it is a
usable one: it says which of the models in that chapter are worth attacking
analytically and which should be left to simulation.

\section{The quadrant argument}
\label{con:sec:quadrant}

The argument here is the crudest in the thesis and, for that reason, the one
whose assumptions are easiest to see. It sorts the bacterium's options into a
two-by-two --- accept or resist phagocytosis, before or after the neutrophil
influx --- and asks which cell of that grid offers the better return. The
claim is that the best cell moves along the diagonal as the influx proceeds,
and that it does so at a threshold which can be written down.

\subsection{Three environments, one ordering}
\label{con:sec:quadrant-ordering}

Simplify the intracellular options to macrophages and neutrophils, setting
aside dendritic cells, eosinophils, basophils and mast cells, which are no
doubt relevant in reality. Macrophages are generalist phagocytes; neutrophils
are specialist killers, and it is mainly neutrophils that are recruited in
numbers once a bacterial infection of this kind is recognised. Treat
replication in each of the three environments --- inside a macrophage, inside
a neutrophil, and extracellular --- as a simple birth--death process with
rates $(\beta_i,\mu_i)$ for $i\in\{M,N,E\}$. The difference in hostility then
sits in the death rate, and

\begin{equation}
\mu_N>\mu_E>\mu_M .
\label{con:eq:ordering}
\end{equation}

The outer inequalities are the substance of the model, and each is an
inference from the strategy itself rather than an independent measurement.
That \yp{} evolved to exploit the macrophage interior at all is the reason
for taking $\mu_M<\mu_E$; if there were no advantage inside, there would be
no reason to go in. That it resists phagocytosis once neutrophils arrive in
force is the reason for taking $\mu_E<\mu_N$; if the neutrophil interior were
the better place, resisting entry to it would be perverse. The argument is
therefore not a derivation of the strategy from first principles. It is a
consistency check: it asks whether one ordering of three death rates can
account for both halves of the observed behaviour at once.

Take the birth rates equal,
\begin{equation}
\beta_M=\beta_E=\beta_N=\beta ,
\label{con:eq:births}
\end{equation}
which is certainly false in detail --- nutrient and energy supply differ
inside a cell and outside it --- but which keeps the comparison to one axis.
What matters to the comparison is the net rate $\beta_i-\mu_i$ in any case,
and by \cref{con:eq:ordering,con:eq:births},
\begin{equation}
\beta_M-\mu_M \;>\; \beta_E-\mu_E \;>\; \beta_N-\mu_N .
\label{con:eq:netordering}
\end{equation}
The macrophage is the best environment available. The bacterium cannot,
however, choose it: the assumption here is that \yp{} can decide whether to
permit phagocytosis but cannot decide which phagocyte performs it.

\begin{remark}[What the argument does not consider]
\label{con:rem:alternatives}
There are other reasons \yp{} might favour macrophage uptake early. The
infected macrophage may serve chiefly as transport to the draining lymph
nodes, where the food supply is better; and the pro-phagocytic phase may be a
lingering hangover from life inside the flea, where the temperature cue for
the second phase is absent. Neither possibility is excluded by anything
below. The argument explores whether a more fundamental strategic advantage
is available in addition, not instead.
\end{remark}

Also folded into a constant here is the mechanism of the killing itself.
Leucocytes kill from a store of granules, and
\mextref{path:sec:intracellular}{\S8.3} shows that when the store is
depletable the effective death rate falls as a cohort spends it --- so
$\mu_M$ and $\mu_N$ are not constants but declining functions of how much
killing has already been done. Taking them fixed is the approximation that
makes the present argument a two-line calculation, and
\cref{con:sec:limits} records it as the place where the argument is weakest.

\subsection{Where the two strategies cross}
\label{con:sec:quadrant-threshold}

Following pathogen recognition the body delivers neutrophils to the site, so
that the composition of the phagocyte population shifts. Write
\begin{equation}
p_M=\frac{\#M}{\#M+\#N},\qquad
p_N=\frac{\#N}{\#M+\#N}=1-p_M ,
\label{con:eq:fractions}
\end{equation}
for the probability that an uptake event is performed by each type, and take
$p_N$ to rise at the influx. A bacterium that permits phagocytosis therefore
draws its environment at random from the two, while one that resists it takes
the extracellular rate with certainty:
\begin{align}
R_{\mathrm{pro}}&=p_M(\beta_M-\mu_M)+p_N(\beta_N-\mu_N),
\label{con:eq:Rpro}\\
R_{\mathrm{anti}}&=\beta_E-\mu_E .
\label{con:eq:Ranti}
\end{align}
Eliminating $p_M$ makes the dependence on the influx explicit,
\begin{equation}
R_{\mathrm{pro}}
=(\beta_M-\mu_M)+p_N\lrbs{(\beta_N-\mu_N)-(\beta_M-\mu_M)} ,
\label{con:eq:Rpro-linear}
\end{equation}
so that $R_{\mathrm{pro}}$ falls linearly in $p_N$, from $\beta_M-\mu_M$ at
$p_N=0$ to $\beta_N-\mu_N$ at $p_N=1$, while $R_{\mathrm{anti}}$ does not
move. The two therefore cross exactly once. Resisting phagocytosis becomes
the better strategy when $R_{\mathrm{anti}}>R_{\mathrm{pro}}$, which by
\cref{con:eq:Rpro-linear} --- and remembering that the bracket is negative,
so that the inequality reverses on division --- occurs when
\begin{equation}
p_N \;>\;
\frac{(\beta_E-\mu_E)-(\beta_M-\mu_M)}{(\beta_N-\mu_N)-(\beta_M-\mu_M)} ,
\label{con:eq:threshold-general}
\end{equation}
and with the birth rates equal this reduces to
\begin{equation}
\boxed{\;p_N \;>\; p_N^{\ast}
:=\frac{\mu_E-\mu_M}{\mu_N-\mu_M}\;}
\label{con:eq:threshold}
\end{equation}

\begin{figure}[!htbp]
\centering
\begin{tikzpicture}[>={Stealth[length=1.9mm]},font=\small]
% axes
\draw[->,inksoft,thick] (0,0) -- (9.3,0)
  node[below left=1pt and -6pt,font=\footnotesize]{$p_N$};
\draw[->,inksoft,thick] (0,0) -- (0,4.3)
  node[above,font=\footnotesize]{net rate};
\node[font=\footnotesize,inksoft] at (0,-0.32) {$0$};
\node[font=\footnotesize,inksoft] at (8.6,-0.32) {$1$};
% R_pro: falls from beta_M-mu_M to beta_N-mu_N
\draw[accent,very thick] (0,3.7) -- (8.6,0.7);
\node[accent,font=\footnotesize,anchor=west] at (0.15,3.98)
  {$R_{\mathrm{pro}}$};
\node[inksoft,font=\scriptsize,anchor=east] at (-0.12,3.7)
  {$\beta_M-\mu_M$};
\node[inksoft,font=\scriptsize,anchor=west] at (8.75,0.7)
  {$\beta_N-\mu_N$};
% R_anti constant
\draw[type2,very thick] (0,2.3) -- (8.6,2.3);
\node[type2,font=\footnotesize,anchor=west] at (7.0,2.62)
  {$R_{\mathrm{anti}}=\beta_E-\mu_E$};
% crossing
\coordinate (X) at (4.01,2.3);
\fill[ink] (X) circle (2.1pt);
\draw[ink,densely dashed] (X) -- (4.01,0);
\node[ink,font=\footnotesize,anchor=north] at (4.01,-0.12)
  {$p_N^{\ast}$};
% regime labels
\node[inksoft,font=\scriptsize,align=center] at (2.0,0.80)
  {accept phagocytosis};
\node[inksoft,font=\scriptsize,align=center] at (6.15,0.80)
  {resist phagocytosis};
% direction of the influx, kept clear of the axis labels
\draw[->,inksoft] (2.6,-1.15) -- (5.9,-1.15);
\node[inksoft,font=\scriptsize,anchor=south] at (4.25,-1.10)
  {neutrophil influx};
\end{tikzpicture}
\caption{The crossing. The expected net replication rate under a
pro-phagocytic strategy falls linearly in the neutrophil fraction $p_N$,
because a larger fraction of uptake events delivers the bacterium into the
more hostile of the two cell types; the anti-phagocytic rate does not depend
on $p_N$ at all. The single crossing is $p_N^{\ast}$ of
\cref{con:eq:threshold}, and it lies strictly between $0$ and $1$ exactly
when the ordering \cref{con:eq:ordering} holds.}
\label{con:fig:quadrant}
\end{figure}

\subsection{What the threshold says}
\label{con:sec:quadrant-reading}

\Cref{con:eq:threshold} has a reading that makes the ordering assumption do
visible work rather than sit in the background. The threshold is the position
of the extracellular death rate on the interval running from the macrophage
rate to the neutrophil rate: if the outside world is a quarter of the way
from macrophage to neutrophil in hostility, the switch comes when a quarter
of the phagocytes on the scene are neutrophils. That reading also shows what
the ordering buys, since $p_N^{\ast}$ lies in $(0,1)$ if and only if
\cref{con:eq:ordering} holds. Were the extracellular environment kinder than
the macrophage interior, the threshold would be negative and resistance would
be right from the first moment; were it harsher than the neutrophil interior,
the threshold would exceed one and the bacterium should accept phagocytosis
however many neutrophils arrive. A biphasic strategy is optimal only in the
intermediate case, and the intermediate case is precisely the ordering
inferred from the strategy. The argument is a consistency check that passes,
and it should be read as no stronger than that.

Two things follow that are not merely restatements. The first is that the
switch is keyed to a composition and not to an elapsed time, so anything that
changes the arriving mixture should move it. Neutrophil depletion should
delay the switch or remove it; accelerated recruitment should bring it
forward. That is the same shape of prediction as the removal budget
$\vartheta^{\ast}$ of \mextref{path:sec:ypestis}{\S8.5.4}, reached from a
different model with different machinery, and the two accounts have not been
reconciled. The second is that $p_N^{\ast}$ is written entirely in death
rates, which are in principle measurable: intracellular survival of \yp{} in
macrophages and in neutrophils, and survival in the extracellular medium, are
three quantities of the same kind, and the argument stands or falls on
whether the middle one lies between the other two.

Finally, the argument as set out is indifferent to \emph{when} the influx
happens. That the bacterium's residence in macrophages may itself delay
neutrophil recruitment, by hiding the infection, is a real effect and would
show up as $t_N$ depending on the strategy chosen. It does not enter here,
because whether uptake is advantageous at a given moment depends on the
mixture present at that moment and not on the history that produced it.

\section{Limitations and open problems}
\label{con:sec:limits}

Nothing in this thesis has been fitted to plague data. The rates the models
need for the intracellular phase of \yp{} do not exist in the literature in
the form required, and the two published parameter sets used for illustration
in \mextref{dist:ch:distribution-theory}{Chapter~5} belong to other
organisms. Everything said about \yp{} is therefore a hypothesis compatible
with the cited intracellular literature, not an inference this work can
check.

Four assumptions carry the closed forms and none is innocuous. Birth is
linear, so the conditional load has no intracellular carrying capacity; the
defence is that the rupture hazard is itself proportional to load and so
regulates it, but the sensitivity of the burst statistics to a logistic birth
rate has not been checked. The rupture hazard is exactly $\delta n$ --- not
$\delta n^{2}$, not a critical load, not a maturation-timed lysis --- and
this is the strongest assumption anywhere in the thesis, because the burst
law and its telescoping proof both depend on it. Release is instantaneous and
complete, which is what makes the burst a single random variable; an assay
that saw graded release would refute the premise before reaching any of the
theory. And the rates are constant in time, which \mextref{var:ch:rates}{Chapter~9}
relaxes only in models that no longer close.

The quadrant argument of \cref{con:sec:quadrant} is the weakest piece of
mathematics in the thesis, and it is kept because it is transparent rather
than because it is strong. It treats $\mu_M$ and $\mu_N$ as constants at
exactly the point where \mextref{path:ch:pathogenesis}{Chapter~8} shows them
not to be: killing from a depletable granule store means the death rate falls
as a cohort spends the store, so the comparison between environments should
depend on cohort size, and in this argument it does not. It also compares
instantaneous rates rather than probabilities of establishment, which is the
quantity that actually decides the fate of a small founding population.

Two structural gaps remain between the chapters. There is no host-to-host
layer: the released particles are counted and then forgotten, so the
population-level consequences of a threshold clearance are not derived
anywhere. And multiplicity of infection, which the single-cell theory handles
cleanly, has not been carried through the population layer, even though the
kernels for it are available and add no fitted parameters.

Of the open problems recorded across the chapters, four are the ones worth
taking first.

\begin{enumerate}[leftmargin=1.8em,itemsep=0.3em]
\item \textbf{Reconcile the two thresholds.} Whether the neutrophil fraction
that crosses $p_N^{\ast}$ in \cref{con:eq:threshold} is the same event as the
toll crossing $\vartheta^{\ast}$ in
\mextref{path:cor:threshold}{Corollary~8.8}. Two accounts of one transition,
in different currencies, is one account too many.
\item \textbf{Two-type killed second moments.} These are what an
eclipse-aware renewal model needs, and they are the principal technical
obstacle between \mextref{p:ch:one-cell-to-population}{Chapter~6} and
\mextref{tt:ch:two-type}{Chapter~7}. The one-type escape was to
differentiate the Riccati equation; the two-type route should be to
differentiate the backward system. It should not be assumed easy.
\item \textbf{Locate \yp{} relative to $L=1$.} This is the question that
would turn the bursting comparison from a criterion into a statement about
the organism, and it needs $\delta$, which needs the exit mechanism.
\item \textbf{The chained process at $\mu>0$.} The joint laws of rupture
times and sizes are known in closed form only when internal extinction is
impossible, and the decomposition fails as soon as a lineage has two ways to
stop rather than one.
\end{enumerate}

\section{Closing}
\label{con:sec:closing}

Set against one switch, what this thesis can offer is two thresholds and a
parameter. The thresholds are $p_N^{\ast}$, the neutrophil fraction at which
being eaten stops paying, and $\vartheta^{\ast}$, the removal budget beyond
which no intracellular strategy stays supercritical; the parameter is
$L=a(1-b)$, which decides whether releasing everything at once is a good bet
or a bad one and changes its answer according to whether the defence facing
the release is a rate or a budget. None of the three is measured. All three
are written in quantities that could be.

That is a narrower return than the question seemed to promise, and the reason
is worth naming, because it is also the pattern that recurred throughout. So
much of the biology turns on arriving in numbers rather than on replicating
quickly. A near-critical population is almost indifferent to its growth rate
and acutely sensitive to how many founders it started with; depletable
defences reward the cohort that spends them together; and the whole
intracellular phase reads, from far enough back, as a device for converting a
slow advantage into a large one at a single instant. The mathematics of that
conversion is what the middle chapters are about, and the burst-size law is
where it is at its most compact: a geometric distribution whose ratio is the
larger root of a quadratic in three rates.

\Mextref{var:ch:rates}{Chapter~9} came from a different appetite, and it
marks plainly which of its parts are results and which are hunches. The
suspicion that drove it --- that tension accumulates and releases in cycles
across scales, and that the release is where new form appears --- remains a
suspicion. What it produced that is defensible is narrower and worth having:
the observation that a process may act on its own circumstances and stay
solvable, so long as what crosses the join is a random time rather than a
bending rate. \yp{} is a good illustration of that principle, since the whole
of its intracellular phase reaches the extracellular one through a single
number, the size of the burst.

What would settle the questions this thesis leaves open is not more theory.
It is three death rates and a rupture rate, measured for one organism in one
cell type: time-lapse imaging of infected macrophages held to rupture, with
the moment of release recorded and the payload counted. The predictions are
sharp enough to be wrong, which is the most that a model of this kind should
be asked to be.
%<<<CHAPTER BODY END>>>
% ===========================================================================

\FloatBarrier
\begingroup
\footnotesize
\bibliographystyle{unsrt}
\bibliography{references}
\endgroup

\end{document}
